\documentclass{beamer}

\usetheme{Madrid}
\usecolortheme{seahorse}
\title{The AKS Algorithm}
\author{Siddharth Narayan Mishra}
\date{\today}
\setbeamertemplate{navigation symbols}{}

\newcommand{\Z}{\mathbb{Z}}
\newcommand{\bigO}{\mathcal{O}}

\begin{document}

\maketitle

%------------------------------------------------
\section{The Breakthrough}

\begin{frame}{The Question}

\begin{center}
\Large
Is primality testing in $\mathbf{P}$?
\end{center}


\begin{itemize}
    \item Fermat: fast but unreliable
    \item Miller-Rabin: fast, probabilistic
\end{itemize}


\begin{block}{AKS (2002)}
\begin{center}
\textbf{PRIMES is in $\mathbf{P}$}
\end{center}
\end{block}

\end{frame}

%------------------------------------------------

\begin{frame}{What AKS Achieves}

\begin{itemize}
    \item Deterministic
    \item Unconditional
    \item Polynomial time
\end{itemize}


\[
\text{Time: } \bigO((\log n)^{c})
\]


\begin{block}{Significance}
First proof that primality testing is efficiently solvable
\end{block}

\end{frame}

%------------------------------------------------
\section{Key Identity}

\begin{frame}{The Core Idea}

\begin{block}{Key Identity}
$n$ is prime iff:
\[
(X+a)^n \equiv X^n + a \pmod{n}
\]
\end{block}


\begin{itemize}
    \item Holds for all $a$ with $\gcd(a,n)=1$
    \item Comes from binomial expansion
\end{itemize}

\end{frame}

%------------------------------------------------

\begin{frame}{Proof ($\Rightarrow$ direction)}

Assume $n$ is prime.

\vspace{0.5em}

\[
(X+a)^n = \sum_{i=0}^{n} \binom{n}{i} a^{n-i} X^i
\]

\vspace{0.5em}

For $0 < i < n$:
\[
\binom{n}{i} = \frac{n!}{i!(n-i)!}
\]
\[
\binom{n}{i} \equiv 0 \pmod{n}
\]

\end{frame}

%------------------------------------------------

\begin{frame}{Proof (continued)}

Thus:
\[
(X+a)^n \equiv X^n + a^n \pmod{n}
\]


By Fermat's Little Theorem:
\[
a^n \equiv a \pmod{n}
\]


\[
\Rightarrow (X+a)^n \equiv X^n + a \pmod{n}
\]

\end{frame}

%------------------------------------------------

\begin{frame}{Proof ($\Leftarrow$ direction)}

Assume:
\[
(X+a)^n \equiv X^n + a \pmod{n}
\]


Suppose $n$ is composite.

Let $q^k \parallel n$ (largest power of $q$ dividing $n$).

\end{frame}

%------------------------------------------------

\begin{frame}{Proof (continued)}

Let $q^k \parallel n$ (i.e., $q^k \mid n$ but $q^{k+1} \nmid n$).

\vspace{0.5em}

Consider:
\[
\binom{n}{q} = \frac{n(n-1)\cdots(n-q+1)}{q!}
\]

\vspace{0.5em}

Rewrite:
\[
\binom{n}{q}
= \frac{q^k (n-1)\cdots(n-q+1)}{q \cdot (q-1)!}
= \frac{q^{k-1}(n-1)\cdots(n-q+1)}{(q-1)!}
\]

\end{frame}
\begin{frame}

\textbf{Key point:}

\begin{itemize}
\item Numerator contains exactly $q^{k-1}$ as the highest power of $q$
\item None of $(n-1), \dots, (n-q+1)$ are divisible by $q$
\item Denominator $(q-1)!$ contains no factor of $q$
\end{itemize}

\vspace{0.5em}

So $\binom{n}{q}$ contains exactly $q^{k-1}$ as a factor.

\vspace{0.5em}

But $n$ contains $q^k$ as a factor.

\[
n \nmid \binom{n}{q}
\quad \Rightarrow \quad
\binom{n}{q} \not\equiv 0 \pmod{n}
\]

\end{frame}

%------------------------------------------------

%------------------------------------------------

\begin{frame}{Example: Prime Case ($n=5$)}

\[
(X+1)^5 = X^5 + 5X^4 + 10X^3 + 10X^2 + 5X + 1
\]


Modulo $5$:
\[
(X+1)^5 \equiv X^5 + 1 \pmod{5}
\]


\begin{block}{Observation}
All middle coefficients vanish modulo $5$
\end{block}

\end{frame}

%------------------------------------------------

\begin{frame}{Pascal Triangle Insight}

Row $5$ of Pascal's Triangle:

\[
\begin{array}{c}
1 \\
1\;1 \\
1\;2\;1 \\
1\;3\;3\;1 \\
1\;4\;6\;4\;1 \\
1\;5\;10\;10\;5\;1
\end{array}
\]

\begin{itemize}
    \item For prime $p$, all interior entries are divisible by $p$
    \item Only the endpoints are $1$
\end{itemize}

\end{frame}

%------------------------------------------------

\begin{frame}{Example: Composite Case ($n=4$)}

\[
(X+1)^4 = X^4 + 4X^3 + 6X^2 + 4X + 1
\]


Modulo $4$:
\[
(X+1)^4 \equiv X^4 + 2X^2 + 1 \pmod{4}
\]


\begin{block}{Observation}
Middle coefficient $6 \equiv 2 \pmod{4}$ does not vanish
\end{block}

\vspace{0.5em}

\[
\Rightarrow \text{Identity fails}
\]

\end{frame}

\begin{frame}{Why Not Use This Directly?}

\begin{itemize}
    \item The identity gives a \textbf{perfect characterization} of primes
\end{itemize}


\[
(X+a)^n = \sum_{i=0}^{n} \binom{n}{i} a^{n-i} X^i
\]


\begin{itemize}
    \item Polynomial has $n+1$ terms
    \item Computing it requires $\Omega(n)$ operations
\end{itemize}


\begin{block}{Problem}
Input size is $\log n$, but runtime is $\Omega(n)$
\end{block}

\vspace{0.5em}

\begin{block}{Conclusion}
Exponential time $\Rightarrow$ not efficient
\end{block}

\end{frame}

%------------------------------------------------
\section{The Trick}

\begin{frame}{Reducing the Cost}

Work modulo:
\[
X^r - 1
\]


\begin{itemize}
    \item $X^r \equiv 1$
    \item Degrees wrap around
\end{itemize}


\begin{block}{Effect}
Polynomial size reduces from $n$ to $r$
\end{block}

\end{frame}

\begin{frame}{Example: Working Modulo $X^r - 1$}

Let $r = 5$. Then:
\[
X^5 \equiv 1 \pmod{X^5 - 1}
\]


\textbf{Example: Reduce } $X^7$

\[
X^7 = X^{5+2} = X^5 \cdot X^2 \equiv 1 \cdot X^2 = X^2
\]


\end{frame}
\begin{frame}
\textbf{Example: Reduce a polynomial}

\[
f(X) = X^7 + 3X^6 + 2X^2
\]

\[
X^7 \equiv X^2,\quad X^6 = X^{5+1} \equiv X
\]

\[
\Rightarrow f(X) \equiv X^2 + 3X + 2X^2 = 3X^2 + 3X
\]


\begin{block}{Key Idea}
Every time degree $\ge r$, replace $X^r$ by $1$ and reduce
\end{block}

\end{frame}

%------------------------------------------------

\begin{frame}{New Test}

Check:
\[
(X+a)^n \equiv X^n + a \pmod{X^r -1,\; n}
\]


\begin{itemize}
    \item Much faster
    \item Works for primes
\end{itemize}


\begin{block}{Issue}
Some composites may now pass
\end{block}

\end{frame}

%------------------------------------------------

\begin{frame}{Fixing False Positives}

\begin{itemize}
    \item Work modulo $X^r - 1$ (compression)
    \item Test multiple values of $a$
\end{itemize}


\begin{block}{Key Condition}
\[
\mathrm{ord}_r(n) > (\log n)^2
\]
\end{block}

\vspace{0.5em}

\begin{itemize}
    \item Ensures composite numbers cannot satisfy identity for many $a$
\end{itemize}

\end{frame}

%------------------------------------------------
\begin{frame}{What is $\mathrm{ord}_r(n)$?}

\begin{block}{Definition}
\[
\mathrm{ord}_r(n) = \min\{k > 0 \mid n^k \equiv 1 \pmod{r}\}
\]
\end{block}


\begin{itemize}
    \item Measures how $n$ behaves modulo $r$
    \item Captures the \textbf{cycle length} of powers of $n$
\end{itemize}


\begin{block}{Example}
\[
2^3 \equiv 1 \pmod{7} \Rightarrow \mathrm{ord}_7(2) = 3
\]
\end{block}

\end{frame}

%------------------------------------------------

\begin{frame}{Why Do We Need Large Order?}

\begin{itemize}
    \item We want:
\[
\mathrm{ord}_r(n) > (\log n)^2
\]
\end{itemize}


\begin{block}{Reason}
Ensures $n$ does not behave ``periodically'' too early modulo $r$
\end{block}


\begin{itemize}
    \item Prevents composite numbers from mimicking primes
    \item Forces large algebraic structure in the proof
\end{itemize}

\end{frame}

%------------------------------------------------

%------------------------------------------------

\begin{frame}{Existence of $r$}

\begin{block}{Lemma}
There exists $r \le (\log n)^5$ such that:
\[
\mathrm{ord}_r(n) > (\log n)^2
\]
\end{block}

\vspace{0.5em}

\textbf{Proof.}

Let:
\[
B = \lceil (\log n)^5 \rceil
\]

Assume for contradiction that for all $r \le B$:
\[
\mathrm{ord}_r(n) \le (\log n)^2
\]

\end{frame}

%------------------------------------------------

\begin{frame}{Key Construction}

Then for every $r \le B$, there exists $k \le (\log n)^2$ such that:
\[
n^k \equiv 1 \pmod{r}
\]

\[
\Rightarrow r \mid (n^k - 1)
\]


Thus every $r \le B$ divides:
\[
N = \prod_{k=1}^{\lfloor (\log n)^2 \rfloor} (n^k - 1)
\]

\end{frame}

%------------------------------------------------

\begin{frame}{Bounding $N$}

Each term:
\[
n^k - 1 < n^k
\]

So:
\[
N < \prod_{k=1}^{(\log n)^2} n^k = n^{\sum_{k=1}^{(\log n)^2} k}
\]

\[
= n^{\frac{(\log n)^2((\log n)^2 + 1)}{2}} \le n^{(\log n)^4}
\]

\[
= 2^{(\log n)^5}
\]

\end{frame}

%------------------------------------------------

\begin{frame}{Contradiction via LCM}

But if every $r \le B$ divides $N$, then:
\[
\mathrm{lcm}(1,2,\dots,B) \mid N
\]


\begin{block}{Fact}
\[
\mathrm{lcm}(1,2,\dots,B) \ge 2^B
\]
\end{block}


Thus:
\[
2^B \le N < 2^{(\log n)^5}
\]

\end{frame}

%------------------------------------------------

\begin{frame}{Final Contradiction}

Recall:
\[
B = (\log n)^5
\]

So:
\[
2^B = 2^{(\log n)^5}
\]

But we showed:
\[
N < 2^{(\log n)^5}
\]


\[
\Rightarrow 2^B \nmid N
\]


\begin{block}{Contradiction}
Not all $r \le B$ can divide $N$
\end{block}

\end{frame}

%------------------------------------------------

\begin{frame}{Conclusion}

\begin{itemize}
    \item There exists some $r \le (\log n)^5$ such that:
\[
\mathrm{ord}_r(n) > (\log n)^2
\]
\end{itemize}

\end{frame}



%------------------------------------------------
\section{The Algorithm}

\begin{frame}{AKS Algorithm}

\begin{enumerate}
    \item Check if $n = a^b$ for some $b > 1$
    \item Find smallest $r$ with $\mathrm{ord}_r(n) > (\log n)^2$
    \item If $\gcd(a,n) > 1$ for some $a \le r$, output composite
    \item If $n \le r$, output prime
    \item For $a = 1$ to $O(\sqrt{r}\log n)$:
    \[
    (X+a)^n \equiv X^n + a \pmod{X^r -1, n}
    \]
    \item Output prime
\end{enumerate}

\end{frame}

%------------------------------------------------

\begin{frame}{Why Each Step Matters}

\begin{itemize}
    \item Step 1: removes perfect powers
    \item Step 2: ensures strong algebraic structure
    \item Step 3: detects small factors quickly
    \item Step 4: handles small cases
    \item Step 5: main polynomial test
\end{itemize}


\begin{block}{Key Point}
All steps together eliminate every possible composite
\end{block}

\end{frame}

%------------------------------------------------

\begin{frame}{Correctness: Setup}

Assume the algorithm reaches the final step and outputs \textsc{Prime}.

\vspace{0.5em}

Then:
\begin{itemize}
    \item $n$ is not a perfect power
    \item $\gcd(a,n)=1$ for all $a \le r$
    \item $\mathrm{ord}_r(n) > (\log n)^2$
    \item For all $1 \le a \le \ell$:
\[
(X+a)^n \equiv X^n + a \pmod{X^r-1,\, n}
\]
\end{itemize}

\vspace{0.5em}

Let $p$ be any prime divisor of $n$.

\end{frame}

%------------------------------------------------

\begin{frame}{Reducing Modulo $p$}

Since $p \mid n$:
\[
(X+a)^n \equiv X^n + a \pmod{X^r-1,\, p}
\]


Also, since $p$ is prime:
\[
(X+a)^p \equiv X^p + a \pmod{p}
\]


Raise both sides to power $n/p$:
\[
(X+a)^n \equiv (X^p + a)^{n/p}
\]

\end{frame}

%------------------------------------------------

\begin{frame}{Key Identity}

Combining:
\[
(X^p + a)^{n/p} \equiv X^n + a \pmod{X^r-1,\, p}
\]


Let $Y = X^p$:
\[
(Y + a)^{n/p} \equiv Y^{n/p} + a
\]


\begin{block}{Observation}
$n/p$ satisfies the same identity as a prime
\end{block}

\end{frame}

%------------------------------------------------

\begin{frame}{Introspective Property}

Define:
\[
f(X)^m \equiv f(X^m) \pmod{X^r-1,\, p}
\]


Then:
\begin{itemize}
    \item $p$ satisfies this for $f(X)=X+a$
    \item $n/p$ also satisfies it
\end{itemize}


\begin{block}{Closure}
Products of such numbers also satisfy this property
\end{block}

\end{frame}

%------------------------------------------------

\begin{frame}{Building a Large Set}

Define:
\[
    I = \left\{\, \left(\frac{n}{p}\right)^{\!i} \cdot p^{j} \;\middle|\; i,j \ge 0 \,\right\}
\]


All elements of $I$ satisfy:
\[
f(X)^m \equiv f(X^m)
\]


Reduce modulo $r$:
\[
G = \{ m \bmod r \mid m \in I \}
\]


\[
|G| > (\log n)^2
\]

\end{frame}

%------------------------------------------------------------
\begin{frame}{$I$ Contains All Powers of $n$}
Recall $n = p \cdot \frac{n}{p}$, so setting $i = j$ gives:
\[
    \left(\frac{n}{p}\right)^{\!i} \cdot p^i \;=\; n^i \;\in\; I \qquad \text{for all } i \ge 0.
\]
\vspace{0.5em}

In particular, $I$ contains:
\[
    n^0,\; n^1,\; n^2,\; \ldots
\]
\end{frame}

%------------------------------------------------------------

\begin{frame}{Defining $G$}
We reduce every element of $I$ modulo $r$:
\[
    G = \{\, m \bmod r \mid m \in I \,\}
\]
\vspace{0.5em}

Think of $G$ as the set of ``fingerprints'' that elements of $I$
leave when divided by $r$.
\vspace{1em}

Since $I$ contains all powers of $n$, $G$ contains in particular:
\[
    n^0 \bmod r,\quad n^1 \bmod r,\quad n^2 \bmod r,\quad \ldots
\]
\vspace{0.5em}

The question is: how many of these are distinct?
\end{frame}

%------------------------------------------------------------

\begin{frame}{How Many Distinct Powers of $n$ Are There mod $r$?}
The powers $n^0, n^1, n^2, \ldots$ eventually repeat modulo $r$.
\vspace{1em}

The number of distinct values among $\{n^i \bmod r\}$ is exactly
$\mathrm{ord}_r(n)$ --- the smallest positive $k$ such that $n^k \equiv 1 \pmod{r}$.
\vspace{1em}

Specifically, the values:
\[
    n^0 \bmod r,\quad n^1 \bmod r,\quad \ldots,\quad n^{\mathrm{ord}_r(n)-1} \bmod r
\]
are all distinct, and after that they repeat.
\end{frame}

%------------------------------------------------------------

\begin{frame}{Conclusion: $|G| > (\log n)^2$}
Since all powers $n^0, n^1, \ldots, n^{\mathrm{ord}_r(n)-1}$ are distinct
modulo $r$, and all of them land in $G$:
\[
    |G| \;\ge\; \mathrm{ord}_r(n)
\]
\vspace{0.5em}

We chose $r$ in Step 2 of AKS precisely so that:
\[
    \mathrm{ord}_r(n) \;>\; (\log n)^2
\]
\vspace{0.5em}

Chaining the two inequalities:
\[
    \boxed{|G| \;\ge\; \mathrm{ord}_r(n) \;>\; (\log n)^2}
\]
\end{frame}

\begin{frame}{Working Modulo a Prime: $\mathbb{F}_r$}
Throughout the proof we work with arithmetic modulo $r$, where $r$ is prime.

This means:
\[
    0, 1, 2, \ldots, r-1
\]
are the only numbers, and we always reduce modulo $r$.

For example, with $r = 5$: \quad $3 + 4 = 2$, \quad $3 \times 4 = 2$.

We call this number system $\mathbb{F}_r$.
\end{frame}

%------------------------------------------------------------

\begin{frame}{Polynomials mod $r$}
We can also consider polynomials whose coefficients live in $\mathbb{F}_r$:
\[
    X^3 + 2X + 4, \qquad X^2 + 1, \qquad \ldots
\]
with all coefficients reduced modulo $r$.

For example, with $r = 5$:
\[
    (2X + 3) + (X + 4) = 3X + 7 = 3X + 2
\]

We call this collection $\mathbb{F}_r[X]$.
\end{frame}

%------------------------------------------------------------

\begin{frame}{Reducing Polynomials Further: Mod $h(X)$}
Just as we reduce numbers modulo $r$, we can reduce \emph{polynomials}
modulo another polynomial $h(X)$.

Every time $h(X)$ appears as a factor, we set it to zero, exactly
like setting $r = 0$ in modular arithmetic.

\begin{example}
If $h(X) = X^2 + 1$ and $r = 5$:
\[
    X^3 \bmod (X^2+1) = X \cdot X^2 = X \cdot (-1) = -X = 4X
\]
since $X^2 \equiv -1 \pmod{h(X)}$.
\end{example}
\end{frame}

\begin{frame}{What Is a Multiplicative Inverse?}
The multiplicative inverse of a number $a$ is the number $b$ such that:
\[
    a \times b = 1
\]

For example, in ordinary arithmetic:
\[
    \text{the inverse of } 3 \text{ is } \frac{1}{3}, \quad
    \text{since } 3 \times \frac{1}{3} = 1.
\]

In $\mathbb{F}_5$, fractions do not exist, but inverses still do:
\[
    \text{the inverse of } 3 \text{ is } 2, \quad
    \text{since } 3 \times 2 = 6 \equiv 1 \pmod{5}.
\]
\end{frame}

%------------------------------------------------------------

\begin{frame}{Irreducible Polynomials}
A polynomial $h(X)$ over $\mathbb{F}_r$ is called \emph{irreducible} if it
cannot be factored into two smaller polynomials, the analogue of a prime number,
but for polynomials.

\begin{example}
Over $\mathbb{F}_2$: \quad $X^2 + X + 1$ is irreducible
(it has no roots: $0^2+0+1=1\ne 0$, $\;1^2+1+1=1\ne 0$).
\end{example}

When we reduce modulo an irreducible $h(X)$, every nonzero polynomial
has a multiplicative inverse, we can always divide by any nonzero element.
We call such a number system a \emph{field}.
\end{frame}

%------------------------------------------------------------

\begin{frame}{The Field $\mathbb{F} = \mathbb{F}_r[X]/(h(X))$}
Putting it all together:
\[
    \mathbb{F} = \mathbb{F}_r[X] / (h(X))
\]
is the collection of all polynomials with coefficients in $\mathbb{F}_r$,
reduced modulo $h(X)$.

It behaves like a number system where:
\begin{itemize}
    \item coefficients are reduced mod $r$,
    \item polynomial degree is kept below $\deg h(X)$,
    \item every nonzero element has an inverse (since $h$ is irreducible).
\end{itemize}

This is the arena in which we will count elements of $\hat{G}$.
\end{frame}

%------------------------------------------------
\begin{frame}{What Is $\hat{G}$?}
Fix an irreducible factor $h(X)$ of $X^r - 1$ over $\mathbb{F}_p$.

\[
    P = \left\{ \prod_{a=0}^{\ell} (X+a)^{e_a} \;\middle|\; e_a \ge 0 , \sum_{a=0}^{\ell}{e_a}<t \right\}
\]
where $\ell = \lfloor \sqrt{\varphi(r)} \log n \rfloor$.

$\hat{G}$ is the set of \emph{distinct} images of polynomials in $P$,
inside the field $\mathbb{F} = \mathbb{F}_p[X]/(h(X))$:
\[
    \hat{G} = \{ f(X) \bmod h(X) \mid f(X) \in P \}
\]
\end{frame}

%------------------------------------------------------------

\begin{frame}{Why Distinct Polynomials Stay Distinct in $\hat{G}$}
\begin{block}{Claim}
If $f(X)$ and $g(X)$ are two distinct polynomials in $P$
of degree $< t = |G|$, then they map to \emph{different} elements in $\hat{G}$.
\end{block}

Suppose not: $f(X) = g(X)$ inside $\mathbb{F}$.

Then for any introspective $m \in I$, applying the introspection property to both sides:
\[
    f(X^m) = [f(X)]^m = [g(X)]^m = g(X^m) \quad \text{in } \mathbb{F}
\]

So $f(X^m) - g(X^m) = 0$ in $\mathbb{F}$, meaning $X^m$ is a root of
the polynomial $Q(Y) = f(Y) - g(Y)$.
\end{frame}

%------------------------------------------------------------
\begin{frame}{Reaching a Contradiction}
As $m$ ranges over all elements of $G$, the values $X^m$ are all
\emph{distinct} in $\mathbb{F}$, because $X$ is a primitive $r$-th
root of unity in $\mathbb{F}$.
\vspace{1em}

So $Q(Y) = f(Y) - g(Y)$ has at least $|G| = t$ distinct roots in $\mathbb{F}$.
\vspace{1em}

But $f$ and $g$ both have degree $< t$, so:
\[
    \deg Q \;<\; t
\]
\vspace{0.5em}

A polynomial of degree $< t$ cannot have $t$ distinct roots in a field.
\vspace{1em}

\begin{block}{Contradiction}
Therefore $f(X) \ne g(X)$ in $\mathbb{F}$. \qed
\end{block}
\end{frame}

%------------------------------------------------------------

\begin{frame}{How Many Distinct Polynomials of Degree $< t$ Are in $P$?}
We need to count products $\prod_{a=0}^{\ell}(X+a)^{e_a}$ with $e_a \ge 0$
and total degree:
\[
    \sum_{a=0}^{\ell} e_a \;\le\; t - 1
\]

This is a standard \emph{stars and bars} count: distributing at most $t-1$
among $\ell + 1$ values.

The number of ways is:
\[
    \binom{t - 1 + \ell + 1}{\ell + 1} = \binom{t + \ell}{\ell + 1}
\]
\end{frame}

%------------------------------------------------------------

\begin{frame}{So $|\hat{G}| \ge \binom{t+\ell}{\ell+1}$}
Since each such polynomial is distinct in $\hat{G}$:
\[
    |\hat{G}| \;\ge\; \binom{t+\ell}{\ell+1}
\]

Now we use two facts:
\begin{itemize}
    \item $t = |G| > \log^2 n$, so $\sqrt{t} > \log n$
    \item $\varphi(r) \ge t$, so $\ell = \lfloor\sqrt{\varphi(r)}\log n\rfloor \ge \lfloor\sqrt{t}\log n\rfloor$
\end{itemize}

Write $m = \lfloor \sqrt{t} \log n \rfloor \ge 2$ (since $n \ge 4$). Then $\ell \ge m$, so:
\[
    \binom{t+\ell}{\ell+1} \;\ge\; \binom{2m+1}{m+1} \;=\; \binom{2m+1}{m}
\]
\end{frame}

%------------------------------------------------------------

\begin{frame}{A Simple Binomial Inequality}
\begin{block}{Claim}
For any integer $m \ge 2$: $\displaystyle\binom{2m+1}{m} > 2^{m+1}$
\end{block}

\begin{proof}
The sum of all binomial coefficients gives $2^{2m+1}$. The middle
term $\binom{2m+1}{m}$ is the largest, and there are $m+1$ terms
from $\binom{2m+1}{m}$ up to $\binom{2m+1}{2m+1}$, so:
\[
    (m+1)\binom{2m+1}{m} \;\ge\; 2^{2m} \implies \binom{2m+1}{m} \;\ge\; \frac{2^{2m}}{m+1} > 2^{m+1}
\]
for $m \ge 2$.
\end{proof}
\end{frame}

%------------------------------------------------------------

\begin{frame}{Completing the Lower Bound}
Applying the inequality with $m = \lfloor\sqrt{t}\log n\rfloor$:
\[
    |\hat{G}| \;\ge\; \binom{2m+1}{m} \;>\; 2^{m+1} \;>\; 2^{\sqrt{t}\log n} \;=\; n^{\sqrt{t}}
\]

Since $t = |G|$:

\begin{block}{Lower Bound}
\[
    |\hat{G}| \;>\; n^{\sqrt{|G|}}
\]
\end{block}
\end{frame}
%------------------------------------------------------------

\begin{frame}{Finding a Polynomial That $\hat{G}$ Satisfies}
Consider the subset of $I$:
\[
    J = \left\{ \left(\frac{n}{p}\right)^i \cdot p^j \;\middle|\; 0 \le i, j \le \lfloor\sqrt{t}\rfloor \right\}
\]
where $t = |G|$.

$J$ has $(\lfloor\sqrt{t}\rfloor + 1)^2 > t$ elements.

But $G$ has only $t$ distinct residues mod $r$.

By pigeonhole, two distinct elements of $J$ must give the
same residue mod $r$:
\[
    m_1 \equiv m_2 \pmod{r}, \qquad m_1 \ne m_2, \quad m_1, m_2 \in J.
\]
\end{frame}

%------------------------------------------------------------

\begin{frame}{What Happens When $n$ Is Not a Power of $p$}
The elements of $J$ have the form $\left(\frac{n}{p}\right)^i p^j$.

If two such elements are equal (not just mod $r$, but as integers):
\[
    \left(\frac{n}{p}\right)^{i_1} p^{j_1} = \left(\frac{n}{p}\right)^{i_2} p^{j_2}
\]
then rearranging gives $n^{i_1 - i_2} = p^{j_2 - j_1 - i_2 + i_1}$,
meaning $n$ is a power of $p$.

So if $n$ is \emph{not} a power of $p$, all elements of $J$ are
\emph{distinct as integers}, giving us genuinely distinct $m_1 \ne m_2$
with $m_1 \equiv m_2 \pmod{r}$.
\end{frame}

%------------------------------------------------------------

\begin{frame}{A Polynomial Satisfied by All of $\hat{G}$}
Since $m_1 \equiv m_2 \pmod{r}$, we have $X^{m_1} = X^{m_2}$ in $\mathbb{F}$
(because $X^r = 1$ in $\mathbb{F}$).

So for any $f(X) \in P$, the introspection property gives:
\[
    [f(X)]^{m_1} = f(X^{m_1}) = f(X^{m_2}) = [f(X)]^{m_2} \quad \text{in } \mathbb{F}
\]

Every element $f(X) \in \hat{G}$ is therefore a root of:
\[
    Q(Y) = Y^{m_1} - Y^{m_2}
\]
\end{frame}

%------------------------------------------------------------

\begin{frame}{Bounding the Degree of $Q(Y)$}
The degree of $Q(Y) = Y^{m_1} - Y^{m_2}$ is $\max(m_1, m_2)$.

Both $m_1$ and $m_2$ lie in $J$, so:
\[
    m_1, m_2 \;\le\; \left(\frac{n}{p}\right)^{\lfloor\sqrt{t}\rfloor} \cdot p^{\lfloor\sqrt{t}\rfloor}
    \;=\; n^{\lfloor\sqrt{t}\rfloor} \;\le\; n^{\sqrt{t}}
\]

Therefore:
\[
    \deg Q \;\le\; n^{\sqrt{t}}
\]

Since a polynomial of degree $d$ has at most $d$ roots in a field,
and all $|\hat{G}|$ elements are roots of $Q$:
\[
    \boxed{|\hat{G}| \;\le\; n^{\sqrt{|G|}}}
\]
\end{frame}

%------------------------------------------------------------

\begin{frame}{The Contradiction}
We now have two bounds on $|\hat{G}|$ simultaneously:

\begin{columns}
\begin{column}{0.48\textwidth}
\begin{block}{Lower Bound}
\[
|\hat{G}| > n^{\sqrt{|G|}}
\]
(always true)
\end{block}
\end{column}
\begin{column}{0.48\textwidth}
\begin{block}{Upper Bound}
\[
|\hat{G}| \le n^{\sqrt{|G|}}
\]
(if $n$ not a power of $p$)
\end{block}
\end{column}
\end{columns}
\vspace{1.5em}

These cannot both hold. So our assumption was wrong:
\[
    n \text{ must be a power of } p.
\]
\end{frame}

%------------------------------------------------------------

\begin{frame}{Conclusion}
We have shown $n = p^k$ for some $k \ge 1$.

But AKS checks at the very first step whether $n$ is a perfect power,
and outputs \textsc{Composite} if so.

Since $n$ passed that check, we must have $k = 1$.

\begin{block}{Final Result}
\[
    n = p \implies n \text{ is prime.} \qquad \blacksquare
\]
\end{block}
\end{frame}
%------------------------------------------------
\begin{frame}{Time Complexity: Overview}
Let $k = \log n$ be the input size (number of bits in $n$).

The AKS algorithm has five steps. We analyze each one and
show every step runs in time polynomial in $k$.

\begin{block}{Goal}
Total time $= \mathcal{O}((\log n)^c)$ for some fixed constant $c$.
\end{block}
\end{frame}

%------------------------------------------------------------

\begin{frame}{Step 1: Perfect Power Check}
Check whether $n = a^b$ for some $a \in \mathbb{N}$, $b > 1$.

We try all exponents $b \le \log n$, and for each $b$, compute
the integer $b$-th root of $n$ using binary search.

Each root computation on a $\log n$-bit number costs $\mathcal{O}((\log n)^2)$,
and we do $\mathcal{O}(\log n)$ of them:
\[
    \text{Step 1 time} = \mathcal{O}((\log n)^3)
\]
\end{frame}

%------------------------------------------------------------

\begin{frame}{Step 2: Finding $r$}
Find the smallest $r$ such that $\mathrm{ord}_r(n) > (\log n)^2$.

By Lemma 4, such an $r$ exists with $r \le \log^5 n$.
We try each candidate $r$ and check whether $n^i \equiv 1 \pmod{r}$
for all $i \le (\log n)^2$, stopping as soon as one fails.

There are $\mathcal{O}(\log^5 n)$ candidates, each requiring
$\mathcal{O}((\log n)^2)$ modular exponentiations costing $\mathcal{O}(\log n)$ each:
\[
    \text{Step 2 time} = \mathcal{O}((\log n)^5 \cdot (\log n)^2) = \mathcal{O}((\log n)^7)
\]
\end{frame}

%------------------------------------------------------------

\begin{frame}{Step 3: GCD Checks}
For each $a$ with $1 \le a \le r$, compute $\gcd(a, n)$.

If any $\gcd$ lies strictly between $1$ and $n$, output \textsc{Composite}.

Each GCD via the Euclidean algorithm costs $\mathcal{O}(\log n)$,
and we do at most $r = \mathcal{O}(\log^5 n)$ of them:
\[
    \text{Step 3 time} = \mathcal{O}((\log n)^5 \cdot \log n) = \mathcal{O}((\log n)^6)
\]
\end{frame}

%------------------------------------------------------------

\begin{frame}{Step 4: Small $n$ Check}
If $n \le r$, output \textsc{Prime} directly.

This is a single comparison and costs $\mathcal{O}(1)$, negligible.

If we reach Step 5, we know $n > r$, so $n$ is large enough
that the polynomial checks are meaningful.
\end{frame}

%------------------------------------------------------------

\begin{frame}{Step 5: The Main Loop}
For each $a$ from $1$ to $\ell = \lfloor\sqrt{\varphi(r)}\log n\rfloor$, verify:
\[
    (X + a)^n \equiv X^n + a \pmod{X^r - 1,\; n}
\]

Since $\varphi(r) \le r = \mathcal{O}(\log^5 n)$:
\[
    \ell = \mathcal{O}\!\left(\sqrt{\log^5 n} \cdot \log n\right) = \mathcal{O}((\log n)^{7/2})
\]

This is the most expensive step, we analyze its cost next.
\end{frame}

%------------------------------------------------------------

\begin{frame}{Cost of Each Polynomial Exponentiation}
Each check requires computing $(X+a)^n \bmod (X^r - 1, n)$.

We use repeated squaring, which takes $\mathcal{O}(\log n)$ multiplications.

Each multiplication involves polynomials of degree $< r$ with
coefficients of size $\log n$, costing $\mathcal{O}(r \log n)$.

So the cost per value of $a$ is:
\[
    \mathcal{O}(\log n) \times \mathcal{O}(r \log n) = \mathcal{O}(r (\log n)^2)
\]
\end{frame}

%------------------------------------------------------------

\begin{frame}{Total Time Complexity}
Putting all steps together, Step 5 dominates:
\[
    \text{Step 5 time} = \ell \cdot \mathcal{O}(r(\log n)^2)
\]
\vspace{0.5em}
Substituting $\ell = \mathcal{O}((\log n)^{7/2})$ and $r = \mathcal{O}((\log n)^5)$:
\[
    = \mathcal{O}\!\left((\log n)^{7/2} \cdot (\log n)^5 \cdot (\log n)^2\right)
    = \mathcal{O}\!\left((\log n)^{21/2}\right)
\]
\vspace{0.5em}

\begin{block}{Conclusion}
AKS runs in time $\mathcal{O}((\log n)^{21/2})$, which is polynomial in
the input size $k = \log n$. Hence \textsc{Primes} $\in$ \textbf{P}.
\end{block}
\end{frame}
%------------------------------------------------
\begin{frame}{Can We Do Better?}
The complexity $\mathcal{O}((\log n)^{21/2})$ is polynomial, but can it be improved?

The bottleneck is Step 2: finding $r$ such that $\mathrm{ord}_r(n) > (\log n)^2$.
A tighter bound on $r$ directly reduces the overall complexity.

Lenstra and Pomerance (2011) showed that a better choice of $r$
and a more careful analysis brings the complexity down to:
\[
    \mathcal{O}((\log n)^6)
\]

Further improvements are possible assuming standard conjectures
in number theory, such as the Artin conjecture on primitive roots.
\end{frame}

%------------------------------------------------------------

\begin{frame}{How Does AKS Compare in Practice?}
\begin{center}
\renewcommand{\arraystretch}{1.6}
\begin{tabular}{lccc}
\hline
\textbf{Algorithm} & \textbf{Deterministic?} & \textbf{Always correct?} & \textbf{Complexity} \\
\hline
Trial division     & Yes & Yes & $\mathcal{O}(\sqrt{n})$ \\
Fermat test        & Yes & No  & $\mathcal{O}((\log n)^2)$ \\
Miller-Rabin      & No  & Yes (w.h.p.) & $\mathcal{O}((\log n)^3)$ \\
AKS                & Yes & Yes & $\mathcal{O}((\log n)^6)$ \\
\hline
\end{tabular}
\end{center}

AKS is the only algorithm that is both deterministic and
provably correct for all inputs, with polynomial runtime.
\end{frame}

%------------------------------------------------------------

\begin{frame}{The Practical Catch}
Despite its theoretical importance, AKS is rarely used in practice.

Miller-Rabin runs much faster and the probability of error can be
made astronomically small by repeating the test a few times.

For a 2048-bit number, Miller-Rabin with 40 rounds has an error
probability of at most $4^{-40} \approx 10^{-24}$, smaller than
the probability of a hardware failure mid-computation.

\begin{block}{The honest summary}
AKS matters not because it is the fastest, but because it
\emph{settles a question}: can primality be decided without randomness?
\end{block}
\end{frame}

%------------------------------------------------------------

\begin{frame}{What AKS Actually Proved}
Before 2002, the best deterministic primality tests either:
\begin{itemize}
    \item ran in superpolynomial time, or
    \item required unproven number-theoretic assumptions.
\end{itemize}

AKS gave the first unconditional, deterministic, polynomial-time
algorithm, settling a question that had been open for decades.

It also demonstrated something deeper: the polynomial identities
that characterize primes (Lemma 1) contain enough information
to \emph{decide} primality efficiently, without any randomness.
\end{frame}

%------------------------------------------------------------

\begin{frame}{Where Does AKS Sit in Complexity Theory?}
Recall the major complexity classes:
\[
    \textbf{P} \subseteq \textbf{BPP} \subseteq \cdots
\]
\vspace{0.5em}

\textbf{BPP} is the class of problems solvable in polynomial time
with a randomized algorithm (bounded error probability).

Miller-Rabin places \textsc{Primes} in \textbf{BPP}.

AKS places \textsc{Primes} unconditionally in \textbf{P}:
\[
    \textsc{Primes} \in \mathbf{P}
\]
\vspace{0.5em}

This is one of the few natural problems known to lie in \textbf{P}
after a long period of sitting only in \textbf{BPP}.
\end{frame}

%------------------------------------------------------------

\begin{frame}{Takeaway}
\begin{center}
\Large
Randomness is not necessary.
\end{center}
\vspace{1.5em}

The AKS algorithm shows that one of the most fundamental questions
in mathematics, \emph{is this number prime?}, can be answered
efficiently, deterministically, and unconditionally.

\begin{block}{Final Conclusion}
\[
    \textsc{Primes} \in \mathbf{P}
\]
Primality testing is not just tractable, it is \emph{provably} tractable,
with no randomness, no unproven assumptions, and no exceptions.
\end{block}
\end{frame}

\end{document}